%============================================================
% Chapter 3: Classic Bit Tricks
% STATUS: PLANNED  (generated by tools/outline.py; replace this file when the chapter is written)
%============================================================
\chapter{Classic Bit Tricks}
\label{ch:tricks}

\begin{tcolorbox}[enhanced, colback=white, colframe=secondaryblue, boxrule=1pt, arc=2mm]
\textbf{\large\textcolor{primaryblue}{The one-line tricks every programmer reuses: powers of two, the lowest and highest set bit, and the next mask with the same popcount.}}
\end{tcolorbox}

\chapterfacts{900}{1500}{3}{6}{30}

\begin{plannedbox}[Planned Contents --- this chapter is not written yet]
\textbf{Part A --- Lessons}
\begin{enumerate}[leftmargin=*, label=\textbf{A\arabic*.}, itemsep=1pt]
    \item Powers of two: $x \,\&\, (x-1) = 0$, rounding up and down to a power of two
    \item The lowest set bit: $x \,\&\, (-x)$, clearing it with $x \,\&\, (x-1)$, counting trailing zeros
    \item The highest set bit and the bit length: \texttt{\_\_lg}, \texttt{\_\_builtin\_clz}, and the all-ones mask below it
    \item Masks of bit ranges: extracting, inserting and clearing bits $[l, r)$
    \item Gosper's hack: the next mask with the same number of set bits
\end{enumerate}
\textbf{Part B --- Patterns} (with the two solved problems of each pattern)
\begin{enumerate}[leftmargin=*, label=\textbf{P\arabic*.}, itemsep=1pt]
    \item \textbf{Powers of two} --- Power of Four {\small\color{darkgray}(LeetCode 342)}, Odd Divisor {\small\color{darkgray}(CF 1475A)}
    \item \textbf{The lowest set bit and carries} --- Minimum Operations to Reduce an Integer to 0 {\small\color{darkgray}(LeetCode 2571)}, Johnny and Another Rating Drop {\small\color{darkgray}(CF 1362C)}
    \item \textbf{Highest and lowest bit positions} --- Smallest Number With All Set Bits {\small\color{darkgray}(LeetCode 3370)}, Check if Bitwise OR Has Trailing Zeros {\small\color{darkgray}(LeetCode 2980)}
\end{enumerate}
\end{plannedbox}
